\documentclass[10pt,a4paper]{amsart}
\usepackage{amsmath,amssymb,amsthm,graphicx}
% Retained minimal preamble for the Stacks Project chapter more-algebra.tex (More on Algebra).
% The chapter file pinned-original-source.tex carries only \input{preamble} on its first line, so no
% preamble text lives inside the chapter itself.  The Stacks Project uses one single shared preamble for
% every chapter; the snapshot stages it identically as sources/stacks-chow/preamble.tex and
% sources/stacks-etale-cohomology/preamble.tex.  The block below is copied mechanically, line for line,
% from sources/stacks-chow/preamble.tex lines 14-18 (the Xy-pic setup) and 154-217 (the native theorem
% declarations together with the mathematical macros).  Nothing was added, removed or redefined; the
% wrapper class and the amsmath/amssymb/amsthm/graphicx base packages are supplied by the assembler.
% Chapter: more-algebra.tex -- More on Algebra.
% For commutative diagrams we use Xy-pic
\usepackage[all]{xy}

% We use 2cell for 2-commutative diagrams.
\xyoption{2cell}

% Theorem environments.
%
\theoremstyle{plain}
\newtheorem{theorem}[subsection]{Theorem}
\newtheorem{proposition}[subsection]{Proposition}
\newtheorem{lemma}[subsection]{Lemma}

\theoremstyle{definition}
\newtheorem{definition}[subsection]{Definition}
\newtheorem{example}[subsection]{Example}
\newtheorem{exercise}[subsection]{Exercise}
\newtheorem{situation}[subsection]{Situation}

\theoremstyle{remark}
\newtheorem{remark}[subsection]{Remark}
\newtheorem{remarks}[subsection]{Remarks}

\numberwithin{equation}{subsection}

% Macros
%
\def\lim{\mathop{\mathrm{lim}}\nolimits}
\def\colim{\mathop{\mathrm{colim}}\nolimits}
\def\Spec{\mathop{\mathrm{Spec}}}
\def\Hom{\mathop{\mathrm{Hom}}\nolimits}
\def\Ext{\mathop{\mathrm{Ext}}\nolimits}
\def\SheafHom{\mathop{\mathcal{H}\!\mathit{om}}\nolimits}
\def\SheafExt{\mathop{\mathcal{E}\!\mathit{xt}}\nolimits}
\def\Sch{\mathit{Sch}}
\def\Mor{\mathop{\mathrm{Mor}}\nolimits}
\def\Ob{\mathop{\mathrm{Ob}}\nolimits}
\def\Sh{\mathop{\mathit{Sh}}\nolimits}
\def\NL{\mathop{N\!L}\nolimits}
\def\CH{\mathop{\mathrm{CH}}\nolimits}
\def\proetale{{pro\text{-}\acute{e}tale}}
\def\etale{{\acute{e}tale}}
\def\QCoh{\mathit{QCoh}}
\def\Ker{\mathop{\mathrm{Ker}}}
\def\Im{\mathop{\mathrm{Im}}}
\def\Coker{\mathop{\mathrm{Coker}}}
\def\Coim{\mathop{\mathrm{Coim}}}

% Boxtimes
%
\DeclareMathSymbol{\boxtimes}{\mathbin}{AMSa}{"02}

%
% Macros for moduli stacks/spaces
%
\def\QCohstack{\mathcal{QC}\!\mathit{oh}}
\def\Cohstack{\mathcal{C}\!\mathit{oh}}
\def\Spacesstack{\mathcal{S}\!\mathit{paces}}
\def\Quotfunctor{\mathrm{Quot}}
\def\Hilbfunctor{\mathrm{Hilb}}
\def\Curvesstack{\mathcal{C}\!\mathit{urves}}
\def\Polarizedstack{\mathcal{P}\!\mathit{olarized}}
\def\Complexesstack{\mathcal{C}\!\mathit{omplexes}}
% \Pic is the operator that assigns to X its picard group, usage \Pic(X)
% \Picardstack_{X/B} denotes the Picard stack of X over B
% \Picardfunctor_{X/B} denotes the Picard functor of X over B
\def\Pic{\mathop{\mathrm{Pic}}\nolimits}
\def\Picardstack{\mathcal{P}\!\mathit{ic}}
\def\Picardfunctor{\mathrm{Pic}}
\def\Deformationcategory{\mathcal{D}\!\mathit{ef}}

% --- appended 2026-10-12 (Chief): the project own definitions of the reference/slogan/history
% display environments, copied mechanically from the shared Stacks preamble lines 8-12
% (\usepackage{verbatim} plus the three \newenvironment lines, which the chapter preamble keeps
% verbatim from the project preamble). Nothing added, removed or redefined: these environments
% carry no mathematical content and the project itself typesets them as comments. Without these
% three lines a counted statement whose own text uses them fails with Environment slogan undefined.
\usepackage{verbatim}
\newenvironment{reference}{\comment}{\endcomment}
\newenvironment{slogan}{\comment}{\endcomment}
\newenvironment{history}{\comment}{\endcomment}
\usepackage{cleveref}
\crefname{originallemma}{original lemma}{original lemmas}
\Crefname{originallemma}{Original lemma}{Original lemmas}
\makeatletter
\newlabel{more-algebra-lemma-syntomic-lci}{{07D3}{source}{}{sourcepointer-1}{}}
\newlabel{more-algebra-lemma-syntomic-lci@cref}{{[originallemma][0][]07D3}{[1][0][]source}{}{}{}}
\makeatother
\crefname{originallemma}{original lemma}{original lemmas}
\Crefname{originallemma}{Original lemma}{Original lemmas}
\makeatletter
\newlabel{algebra-definition-lci}{{00SL}{source}{}{sourcepointer-2}{}}
\newlabel{algebra-definition-lci@cref}{{[originallemma][0][]00SL}{[1][0][]source}{}{}{}}
\makeatother
\crefname{originallemma}{original lemma}{original lemmas}
\Crefname{originallemma}{Original lemma}{Original lemmas}
\makeatletter
\newlabel{algebra-lemma-syntomic}{{00SY}{source}{}{sourcepointer-3}{}}
\newlabel{algebra-lemma-syntomic@cref}{{[originallemma][0][]00SY}{[1][0][]source}{}{}{}}
\makeatother
\crefname{originallemma}{original lemma}{original lemmas}
\Crefname{originallemma}{Original lemma}{Original lemmas}
\makeatletter
\newlabel{lemma-lci-local}{{07D1}{source}{}{sourcepointer-4}{}}
\newlabel{lemma-lci-local@cref}{{[originallemma][0][]07D1}{[1][0][]source}{}{}{}}
\makeatother
\crefname{originallemma}{original lemma}{original lemmas}
\Crefname{originallemma}{Original lemma}{Original lemmas}
\makeatletter
\newlabel{lemma-relative-global-complete-intersection-koszul}{{07D2}{source}{}{sourcepointer-5}{}}
\newlabel{lemma-relative-global-complete-intersection-koszul@cref}{{[originallemma][0][]07D2}{[1][0][]source}{}{}{}}
\makeatother
\crefname{originallemma}{original lemma}{original lemmas}
\Crefname{originallemma}{Original lemma}{Original lemmas}
\makeatletter
\newlabel{algebra-definition-lci-field}{{00S9}{source}{}{sourcepointer-6}{}}
\newlabel{algebra-definition-lci-field@cref}{{[originallemma][0][]00S9}{[1][0][]source}{}{}{}}
\makeatother
\crefname{originallemma}{original lemma}{original lemmas}
\Crefname{originallemma}{Original lemma}{Original lemmas}
\makeatletter
\newlabel{definition-regular-ideal}{{07CV}{source}{}{sourcepointer-7}{}}
\newlabel{definition-regular-ideal@cref}{{[originallemma][0][]07CV}{[1][0][]source}{}{}{}}
\makeatother
\crefname{originallemma}{original lemma}{original lemmas}
\Crefname{originallemma}{Original lemma}{Original lemmas}
\makeatletter
\newlabel{lemma-relative-regular-immersion-algebra}{{0CEP}{source}{}{sourcepointer-8}{}}
\newlabel{lemma-relative-regular-immersion-algebra@cref}{{[originallemma][0][]0CEP}{[1][0][]source}{}{}{}}
\makeatother
\crefname{originallemma}{original lemma}{original lemmas}
\Crefname{originallemma}{Original lemma}{Original lemmas}
\makeatletter
\newlabel{lemma-noetherian-finite-all-equivalent}{{09CC}{source}{}{sourcepointer-9}{}}
\newlabel{lemma-noetherian-finite-all-equivalent@cref}{{[originallemma][0][]09CC}{[1][0][]source}{}{}{}}
\makeatother
\crefname{originallemma}{original lemma}{original lemmas}
\Crefname{originallemma}{Original lemma}{Original lemmas}
\makeatletter
\newlabel{algebra-lemma-lci}{{00SC}{source}{}{sourcepointer-10}{}}
\newlabel{algebra-lemma-lci@cref}{{[originallemma][0][]00SC}{[1][0][]source}{}{}{}}
\makeatother
\title{Syntomic ring maps are exactly the flat local complete intersections}
\author{The Stacks Project Authors}
\date{stacks-project:more-algebra.tex:local-snapshot-20261012; GNU FDL 1.2 or later}
\begin{document}
\maketitle
\setcounter{subsection}{1}
\noindent\textit{Editorial note.} All delivered mathematical text is the author's own, copied from the pinned Stacks Project chapter file more-algebra.tex; no statement or proof step was rewritten, reordered or removed.\par\medskip
\noindent\textit{Editorial note.} The chapter file pinned-original-source.tex carries only \textbackslash{}input\{preamble\} on its first line, so no preamble text lives inside the chapter. The retained file sources/stacks-more-algebra/retained-preamble.tex is a mechanical, line-for-line copy of the project's single shared preamble as staged in the snapshot (Xy-pic setup plus theorem declarations and mathematical macros). Nothing was added, removed or redefined, and the excerpt serves only as preamble, never as a span coordinate.\par\medskip
\noindent\textit{Editorial note.} Both counted spans start at the \textbackslash{}begin line of their environment, so the statement and the author's proof are delivered whole; surrounding prose of the section is deliberately outside the counted spans and is carried as standard prerequisites instead.\par\medskip
\noindent\textit{Editorial note.} The citation fields (source\_id, source\_version, authors, source\_url, license.name, source\_title) are identical to those of every other card this lane has produced from this chapter, so the batch selector treats them as one source.\par\medskip
\noindent\textit{Allowed standard prerequisites (editorial).}\begin{itemize}
\item Characterisations of syntomic maps and the reduction to a global complete intersection over a field. Algebra, Lemma \textbackslash{}ref\{algebra-lemma-syntomic\}, Algebra, Definition \textbackslash{}ref\{algebra-definition-lci\}, and Algebra, Definition \textbackslash{}ref\{algebra-definition-lci-field\}, cited at proof lines 8417-8418 and 8427. Reduce the verification to the case of a base change to a field.
\item A relative global complete intersection is a local complete intersection. Lemma \textbackslash{}ref\{lemma-relative-global-complete-intersection-koszul\} and Lemma \textbackslash{}ref\{lemma-lci-local\}, cited at proof lines 8417-8420. Gives one direction of the equivalence.
\item Flat base change turns a Koszul regular sequence into an H1-regular sequence generating the base changed ideal. Lemma \textbackslash{}ref\{lemma-relative-regular-immersion-algebra\}, cited at proof lines 8441-8442. Transports the local presentation of the ideal through the base change to the residue field.
\item A finite module cut out by an H1-regular sequence at a prime is cut out by a regular sequence there. Lemma \textbackslash{}ref\{lemma-noetherian-finite-all-equivalent\} and Algebra, Lemma \textbackslash{}ref\{algebra-lemma-lci\}, cited at proof lines 8445-8448. Upgrade the base changed presentation to a global complete intersection over the field.
\end{itemize}\medskip
\ifdefined\hypertarget\hypertarget{sourcepointer-1}{}\fi
\noindent\textit{Source reference.} Original lemma 07D3: focused-seeds/AGP/sources/stacks-snapshot-20261011/tags.txt, source lines 9398--9398. The lemma's own canonical tag, quoted so that readers can resolve the label inside the delivered text; it points at the counted statement itself and carries no proof content.\par\medskip
\ifdefined\hypertarget\hypertarget{sourcepointer-2}{}\fi
\noindent\textit{Source reference.} Original lemma 00SL: focused-seeds/AGP/sources/stacks-snapshot-20261011/tags.txt, source lines 1043--1043. Navigation pointer to a statement the counted proof refers to but does not contain; nothing load-bearing is delegated to it.\par\medskip
\ifdefined\hypertarget\hypertarget{sourcepointer-3}{}\fi
\noindent\textit{Source reference.} Original lemma 00SY: focused-seeds/AGP/sources/stacks-snapshot-20261011/tags.txt, source lines 1056--1056. Navigation pointer to a statement the counted proof refers to but does not contain; nothing load-bearing is delegated to it.\par\medskip
\ifdefined\hypertarget\hypertarget{sourcepointer-4}{}\fi
\noindent\textit{Source reference.} Original lemma 07D1: focused-seeds/AGP/sources/stacks-snapshot-20261011/tags.txt, source lines 9396--9396. Navigation pointer to a statement the counted proof refers to but does not contain; nothing load-bearing is delegated to it.\par\medskip
\ifdefined\hypertarget\hypertarget{sourcepointer-5}{}\fi
\noindent\textit{Source reference.} Original lemma 07D2: focused-seeds/AGP/sources/stacks-snapshot-20261011/tags.txt, source lines 9397--9397. Navigation pointer to a statement the counted proof refers to but does not contain; nothing load-bearing is delegated to it.\par\medskip
\ifdefined\hypertarget\hypertarget{sourcepointer-6}{}\fi
\noindent\textit{Source reference.} Original lemma 00S9: focused-seeds/AGP/sources/stacks-snapshot-20261011/tags.txt, source lines 1031--1031. Navigation pointer to a statement the counted proof refers to but does not contain; nothing load-bearing is delegated to it.\par\medskip
\ifdefined\hypertarget\hypertarget{sourcepointer-7}{}\fi
\noindent\textit{Source reference.} Original lemma 07CV: focused-seeds/AGP/sources/stacks-snapshot-20261011/tags.txt, source lines 9390--9390. Navigation pointer to a statement the counted proof refers to but does not contain; nothing load-bearing is delegated to it.\par\medskip
\ifdefined\hypertarget\hypertarget{sourcepointer-8}{}\fi
\noindent\textit{Source reference.} Original lemma 0CEP: focused-seeds/AGP/sources/stacks-snapshot-20261011/tags.txt, source lines 15584--15584. Navigation pointer to a statement the counted proof refers to but does not contain; nothing load-bearing is delegated to it.\par\medskip
\ifdefined\hypertarget\hypertarget{sourcepointer-9}{}\fi
\noindent\textit{Source reference.} Original lemma 09CC: focused-seeds/AGP/sources/stacks-snapshot-20261011/tags.txt, source lines 11825--11825. Navigation pointer to a statement the counted proof refers to but does not contain; nothing load-bearing is delegated to it.\par\medskip
\ifdefined\hypertarget\hypertarget{sourcepointer-10}{}\fi
\noindent\textit{Source reference.} Original lemma 00SC: focused-seeds/AGP/sources/stacks-snapshot-20261011/tags.txt, source lines 1034--1034. Navigation pointer to a statement the counted proof refers to but does not contain; nothing load-bearing is delegated to it.\par\medskip
\begin{lemma}
\label{lemma-syntomic-lci}
Let $R \to S$ be a ring map. The following are equivalent
\begin{enumerate}
\item $R \to S$ is syntomic
(Algebra, Definition \ref{algebra-definition-lci}), and
\item $R \to S$ is flat and a local complete intersection.
\end{enumerate}
\end{lemma}

\begin{proof}
Assume (1). Then $R \to S$ is flat by definition.
By Algebra, Lemma \ref{algebra-lemma-syntomic} and
Lemma \ref{lemma-lci-local} we see that it suffices to
show a relative global complete intersection is
a local complete intersection homomorphism which is
Lemma \ref{lemma-relative-global-complete-intersection-koszul}.

\medskip\noindent
Assume (2). A local complete intersection is of finite presentation
because a Koszul-regular ideal is finitely generated.
Let $R \to k$ be a map to a field. It suffices to show that
$S' = S \otimes_R k$ is a local complete intersection over $k$, see
Algebra, Definition \ref{algebra-definition-lci-field}.
Choose a prime $\mathfrak q' \subset S'$.
Write $S = R[x_1, \ldots, x_n]/I$.
Then $S' = k[x_1, \ldots, x_n]/I'$ where
$I' \subset k[x_1, \ldots, x_n]$ is the image of $I$.
Let $\mathfrak p' \subset k[x_1, \ldots, x_n]$,
$\mathfrak q \subset S$,
and $\mathfrak p \subset R[x_1, \ldots, x_n]$
be the corresponding primes.
By Definition \ref{definition-regular-ideal}
exists an $g \in R[x_1, \ldots, x_n]$, $g \not \in \mathfrak p$
and $f_1, \ldots, f_r \in R[x_1, \ldots, x_n]_g$ which form
a Koszul-regular sequence generating $I_g$.
Since $S$ and hence $S_g$ is flat over $R$
we see that the images $f'_1, \ldots, f'_r$ in
$k[x_1, \ldots, x_n]_g$ form a $H_1$-regular sequence
generating $I'_g$, see
Lemma \ref{lemma-relative-regular-immersion-algebra}.
Thus $f'_1, \ldots, f'_r$ map to a regular sequence in
$k[x_1, \ldots, x_n]_{\mathfrak p'}$ generating
$I'_{\mathfrak p'}$ by Lemma \ref{lemma-noetherian-finite-all-equivalent}.
Applying Algebra, Lemma \ref{algebra-lemma-lci}
we conclude $S'_{gg'}$ for some $g' \in S$, $g' \not \in \mathfrak q'$
is a global complete intersection over $k$ as desired.
\end{proof}

\end{document}
